% TEMPLATE-MILC-v3.tex - plantilla maestra UNICA de las guias MILC v3.
%
% La usan las tres familias de guias, cada una con su driver:
%   · Grados 6-11  -> scripts/build-guias-g11.py
%   · Bebras       -> scripts/build-guias-bebras.py
%   · Semillero    -> scripts/build-guias-semillero.py
%
% Antes habia tres copias casi identicas (~400 lineas iguales, 6 distintas):
% cada mejora habia que aplicarla tres veces, y en la practica no se hacia
% (el motor de recursos vivia solo en dos de ellas). Lo que varia entre
% familias esta parametrizado en 6 placeholders que cada driver llena
% (se nombran aqui SIN su sintaxis real: el builder reemplaza por texto plano,
% tambien dentro de los comentarios, y un valor multilinea rompe el preambulo):
%
%   HEADER_IZQ        encabezado izquierdo de pagina
%   HEADER_DER        encabezado derecho de pagina
%   PORTADA_ETIQUETA  rotulo sobre el numero grande (GRADO/MOMENTO/MODULO)
%   PORTADA_SERIE     linea de serie de la portada
%   PORTADA_ID        identificador de la guia en la portada
%   PORTADA_CIRCULO   el circulito de grado (°), o vacio si no aplica
%
% El resto de claves las documenta content/guias/_SCHEMA.md. El script valida
% que no queden placeholders sin reemplazar antes de escribir el archivo final.

\documentclass[11pt,letterpaper]{article}
\usepackage[margin=1.75cm]{geometry}
\usepackage[table]{xcolor}
\usepackage{fontspec}
\usepackage[spanish]{babel}
\usepackage{tikz}
% Librerías TikZ para los diagramas de las guías (bloque `recursos:`):
%  · babel  → babel en español vuelve activos caracteres como «>», y TikZ los usa
%             en sus opciones (>=stealth, ->). Sin esta librería, cualquier
%             diagrama con flechas revienta con «\language@active@arg> has an
%             extra }». Debe ir ANTES de que se dibuje nada.
%  · positioning        → sintaxis «below left=2pt and 3pt».
%  · decorations.pathreplacing → llaves ({brace}) para señalar rangos.
\usetikzlibrary{babel,positioning,decorations.pathreplacing}
\usepackage{graphicx}
% Circuitos: el trabajo de electrónica se hace hoy con micro:bit y sus sensores
% integrados (MakeCode, sin cableado), así que no se necesitan esquemáticos.
% Si algún día entran montajes con protoboard, basta descomentar esta línea:
% los diagramas .tex/.tikz declarados en `recursos:` ya se incrustan solos.
% \usepackage{circuitikz}
\usepackage[most]{tcolorbox}
\usepackage{tabularx}
\usepackage{array}
\usepackage{fancyhdr}
\usepackage{titlesec}
\usepackage{ragged2e}
\usepackage{enumitem}
\usepackage{microtype}

% Helvetica Neue no trae la flecha U+2192, asi que cada «→» escrito literalmente
% en el YAML se compilaba a NADA, en silencio (462 flechas en 61 guias: rutas del
% tipo «paleta → tipografia → lockup» salian rotas). Mapeamos el caracter a su
% version matematica, que si tiene glifo. Asi el autor puede seguir escribiendo
% «→» comodamente en el YAML.
\usepackage{newunicodechar}
\newunicodechar{→}{\ensuremath{\rightarrow}}

\setmainfont{Helvetica Neue}
\setsansfont{Helvetica Neue}
\setlength{\parindent}{0pt}
\setlength{\parskip}{6pt}
\hyphenpenalty=200
\exhyphenpenalty=200
\pretolerance=400
\tolerance=2000
\setlength{\emergencystretch}{1em}
\renewcommand{\arraystretch}{1.25}
\setlist{leftmargin=5mm,itemsep=1mm,topsep=1mm}
\newcolumntype{Y}{>{\RaggedRight\arraybackslash}X}

\definecolor{milcMagenta}{HTML}{E6007E}
\definecolor{milcVino}{HTML}{5A0038}
\definecolor{milcTurquesa}{HTML}{0093A5}
\definecolor{milcVerde}{HTML}{58A923}
\definecolor{milcMostaza}{HTML}{E5B400}
\definecolor{milcOcre}{HTML}{B86600}
\definecolor{milcNegro}{HTML}{241816}
\definecolor{milcGris}{HTML}{F7F7F4}
\definecolor{milcLinea}{HTML}{E8E2DD}
\definecolor{milcGradePrimary}{HTML}{0066FF}
\definecolor{milcGradeDark}{HTML}{003D99}
\definecolor{milcGradeSoft}{HTML}{D6E8FF}
\definecolor{milcGradeAccent}{HTML}{CFE7FF}
\definecolor{milcGradeText}{HTML}{FFFFFF}
\color{milcNegro}

\pagestyle{fancy}
\fancyhf{}
\lhead{\small Tecnología e Informática · Grado 11}
\rhead{\small Guía 16 · MILC v3}
\cfoot{\small\thepage}
\renewcommand{\headrulewidth}{0pt}

\newtcolorbox{titlebox}[1]{
  enhanced,colback=#1,colframe=#1,arc=7mm,boxrule=0pt,
  left=7mm,right=7mm,top=3mm,bottom=3mm,
  before skip=8pt,after skip=8pt,
  fontupper=\sffamily\bfseries\Large\color{white}
}

\newtcolorbox{softbox}[3]{
  enhanced,colback=#2,colframe=#1,borderline west={4pt}{0pt}{#1},
  arc=2mm,boxrule=.4pt,left=4mm,right=4mm,top=3mm,bottom=3mm,
  before skip=6pt,after skip=8pt,title={#3},coltitle=white,
  colbacktitle=#1,fonttitle=\sffamily\bfseries,fontupper=\RaggedRight,
  attach boxed title to top left={xshift=3mm,yshift=-2mm},
  boxed title style={arc=2mm, boxrule=0pt}
}

\newtcolorbox{drawbox}[2]{
  enhanced,colback=white,colframe=#1,arc=4mm,boxrule=1pt,
  left=5mm,right=5mm,top=4mm,bottom=4mm,before skip=8pt,after skip=8pt,
  title={#2},coltitle=white,colbacktitle=#1,fonttitle=\sffamily\bfseries,
  fontupper=\RaggedRight,
  attach boxed title to top left={xshift=4mm,yshift=-2mm},
  boxed title style={arc=3mm, boxrule=0pt}
}

\newtcolorbox{infoband}[1]{
  enhanced,colback=#1!8,colframe=#1!50,arc=2mm,boxrule=.5pt,
  left=4mm,right=4mm,top=2mm,bottom=2mm,before skip=4pt,after skip=4pt,
  fontupper=\small\RaggedRight
}

% Puente narrativo entre secciones (contrato editorial Conéctate).
% Da continuidad: "Acabas de X. Ahora vas a Y porque Z."
% Visualmente: banda izquierda en color del grado, texto en itálica pequeña.
\newtcolorbox{puentebox}{
  enhanced,colback=milcGradeSoft,colframe=milcGradeDark,
  borderline west={3pt}{0pt}{milcGradeDark},
  arc=0mm,boxrule=0pt,left=5mm,right=4mm,top=2.5mm,bottom=2.5mm,
  before skip=4pt,after skip=8pt,
  fontupper=\itshape\small\color{milcNegro}\RaggedRight
}
% La flecha va en modo matematico a proposito: Helvetica Neue NO tiene el glifo
% U+2192, asi que \textrightarrow se compilaba a NADA («Missing character: There
% is no → in font Helvetica Neue») y el puente salia sin flecha. En modo
% matematico la flecha viene de la fuente matematica y si se dibuja.
\newcommand{\puente}[1]{%
  \begin{puentebox}%
    \textcolor{milcGradeDark}{$\rightarrow$}\hspace{2mm}#1%
  \end{puentebox}%
}

\newcommand{\linead}[1]{\noindent\color{milcLinea}\rule{#1}{0.5pt}\color{milcNegro}}
\newcommand{\respuesta}[1]{\par\vspace{2mm}\foreach \n in {1,...,#1}{\noindent\color{milcLinea}\rule{\linewidth}{0.5pt}\par\vspace{5mm}}\color{milcNegro}}
\newcommand{\checkbox}{\textcolor{milcGradeDark}{\fbox{\phantom{x}}}\hspace{5pt}}

\newcommand{\phasecard}[4]{
\begin{tcolorbox}[enhanced,colback=#3,colframe=#2,arc=3mm,boxrule=.5pt,height=3.05cm,valign=center,halign=center,left=1.5mm,right=1.5mm,top=2mm,bottom=2mm]
{\sffamily\bfseries\fontsize{22}{24}\selectfont\textcolor{#2}{#1}}\par
{\sffamily\bfseries\small #4}
\end{tcolorbox}
}

% ── Recursos visuales de la guía ──────────────────────────────────────
% Declarados en el bloque `recursos:` del YAML y emitidos por el builder.
% \guiaFigura   → imagen rasterizada o PDF (foto, carta celeste, montaje).
% \guiaDiagrama → fuente TikZ/circuitikz (.tex) incrustada como vector:
%                 es la vía para circuitos y esquemáticos editables.
% #1 = ruta relativa al .tex · #2 = pie (puede ir vacío)
\newcommand{\guiaFigura}[2]{%
\begin{center}
\includegraphics[width=0.88\linewidth,height=0.38\textheight,keepaspectratio]{#1}\\[1.5mm]
{\footnotesize\itshape\textcolor{milcNegro!75}{#2}}
\end{center}
}

\newcommand{\guiaDiagrama}[2]{%
\begin{center}
\input{#1}\\[1.5mm]
{\footnotesize\itshape\textcolor{milcNegro!75}{#2}}
\end{center}
}

\begin{document}
\thispagestyle{empty}
\begin{tikzpicture}[remember picture,overlay]
  \fill[white] (current page.south west) rectangle (current page.north east);
  \path[fill=milcGradePrimary,rounded corners=16mm]
    ([xshift=.55cm,yshift=-.55cm]current page.north west) rectangle
    ([xshift=-.55cm,yshift=3.65cm]current page.south east);

  \node[anchor=north west,text=milcGradeText] at ([xshift=2.0cm,yshift=-1.85cm]current page.north west)
    {\sffamily\bfseries\fontsize{14}{16}\selectfont GRADO};
  \node[anchor=north west,text=milcGradeText,opacity=.82] at ([xshift=2.0cm,yshift=-2.75cm]current page.north west)
    {\sffamily\bfseries\fontsize{9}{11}\selectfont SERIE GUÍAS · TECNOLOGÍA E INFORMÁTICA};

  \begin{scope}[shift={([xshift=-3.05cm,yshift=-2.55cm]current page.north east)}]
    \fill[white,opacity=.80] (-.62,-.52) rectangle (.62,.52);
    \draw[milcGradeText,line width=1pt] (-.62,-.52) rectangle (.62,.52);
    \fill[milcGradeDark] (-.42,-.38) rectangle (-.22,.06);
    \fill[milcGradeAccent] (-.10,-.38) rectangle (.10,.24);
    \fill[milcGradePrimary!60!black] (.22,-.38) rectangle (.42,.38);
  \end{scope}

  \node[anchor=west,text=milcGradeText] at ([xshift=1.9cm,yshift=-12.85cm]current page.north west)
    {\sffamily\bfseries\fontsize{124}{124}\selectfont 11};
  % El circulito de grado (°) solo aplica a las guias de grado («11°»); en
  % Bebras y Semillero el numero grande es un momento/modulo, no un grado.
  \draw[milcGradeText,line width=1.7pt]
    ([xshift=4.52cm,yshift=-11.68cm]current page.north west) circle (.13cm);

  \node[anchor=west,text=milcGradeText,text width=8.7cm,align=left] at ([xshift=10.35cm,yshift=-11.6cm]current page.north west)
    {
     {\sffamily\bfseries\fontsize{19}{22}\selectfont IA en mensajería ---\\clasificar\\con criterio}\\[4mm]
     {\sffamily\bfseries\fontsize{23}{26}\selectfont Guía 16-11-TIC}\\[2mm]
     {\sffamily\fontsize{13}{16}\selectfont Período 2 · Automatización y procesos}\\[5mm]
     {\sffamily\bfseries\fontsize{11}{13}\selectfont Institución Educativa Sor María Juliana}\\[1mm]
     {\sffamily\fontsize{10}{12}\selectfont Autor: Dr. Álvaro Cárdenas Orozco}
    };

  \path[fill=milcGradeSoft,rounded corners=7mm]
    ([xshift=1.35cm,yshift=.85cm]current page.south west) rectangle
    ([xshift=-1.35cm,yshift=4.25cm]current page.south east);
  \draw[milcGradeDark,line width=.65pt,rounded corners=7mm]
    ([xshift=1.35cm,yshift=.85cm]current page.south west) rectangle
    ([xshift=-1.35cm,yshift=4.25cm]current page.south east);
  \node[anchor=center,text=milcNegro,text width=17.0cm,align=center] at ([yshift=2.55cm]current page.south)
    {\sffamily\fontsize{10}{12}\selectfont Metodología MILC · Apertura ancestral · Pensamiento computacional · Triángulo Dussel-Estoicismo-Floridi\\
    \textcolor{milcGradeDark}{\bfseries Producto:} 1 sistema de clasificación de mensajes (mínimo 20 mensajes reales o simulados) con prompt afinado que separa los mensajes en 3 categorías (urgente / útil / ignorar) con justificación de cada decisión y métrica de aciertos del modelo};
\end{tikzpicture}
\mbox{}
\newpage

\section*{Apertura: IA en mensajería --- clasificar y responder con criterio}

\begin{softbox}{milcGradeDark}{milcGradeSoft}{Saber ancestral}
El \textbf{intermediario respetado de la plaza} de los pueblos del Valle cumplía una función social precisa: filtraba qué noticia se contaba, a quién, en qué momento. Cuando llegaba una novedad al pueblo, no se difundía cruda: pasaba por el \emph{filtro del vecino respetado} (el sacristán, la matrona, el tendero antiguo) que decidía con criterio si era \emph{urgente} (\emph{``hay que avisar al alcalde''}), \emph{útil} (\emph{``el domingo lo cuento en misa''}) o \emph{ignorable} (\emph{``rumor sin verificar''}). La \textbf{abuela jefa de la casa} hacía lo mismo con los mensajes para los niños: filtraba qué noticias del barrio se compartían en la mesa y cuáles se guardaban. Filtrar mensajes con criterio es oficio antiguo. La \textbf{IA en mensajería} es ese filtro --- pero solo es buena si el humano define las reglas con ética.
\end{softbox}

\begin{softbox}{milcTurquesa}{milcGris}{Saber contemporáneo}
La \textbf{clasificación de mensajes con IA} usa modelos de lenguaje (ChatGPT, Claude, Gemini, locales como Llama) para etiquetar automáticamente mensajes según categorías que tú defines. El insumo clave no es el modelo: es el \textbf{prompt} (la instrucción que le das al modelo). Un prompt afinado contiene: (1) \textbf{rol del modelo} (\emph{``actúa como asistente de atención al cliente''}); (2) \textbf{categorías exhaustivas y excluyentes} (\emph{urgente / útil / ignorar}); (3) \textbf{criterios precisos} para cada categoría con ejemplos; (4) \textbf{formato de salida} estructurado (JSON, etiqueta + justificación); (5) \textbf{borde de duda} (\emph{``si dudas, marca útil y deja nota''}). La regla del oficio: \emph{el modelo es solo tan bueno como el prompt; el prompt es solo tan ético como la cabeza que lo escribe}.
\end{softbox}

\begin{softbox}{milcMostaza}{milcGris}{Pregunta-puente}
\textit{¿Cómo sabía la abuela qué noticia merecía atención inmediata, cuál podía esperar y cuál no valía la pena contar? ¿Y qué pierde un emprendedor novato cuando deja a la IA clasificar sus mensajes sin haber definido criterios éticos y reglas claras?}
\end{softbox}

\begin{softbox}{milcVerde}{milcGris}{Saber y saber hacer}
Al terminar podrás: (1) \textbf{identificar} los 5 elementos de un prompt afinado para clasificación; (2) \textbf{analizar} 10 mensajes reales y clasificarlos a mano antes de pasarlos al modelo (línea base humana); (3) \textbf{aplicar} el prompt al lote de mensajes y comparar el resultado del modelo con tu clasificación manual; (4) \textbf{evaluar} la tasa de acierto y los tipos de error del modelo, y decidir si el sistema sirve para producción o necesita afinar el prompt.
\end{softbox}



\small
\begin{tabularx}{\textwidth}{>{\bfseries}p{3.0cm}Y}
\rowcolor{milcGradePrimary!12}Autor & Dr. Álvaro Cárdenas Orozco\\
DBA & Diseño y mejora de procesos digitales con criterio ético (MEN, T\&I)\\
Referentes & Pensamiento computacional · Lean / Kaizen · Floridi (ética informacional)\\
Producto final & 1 sistema de clasificación de mensajes (mínimo 20 mensajes reales o simulados) con prompt afinado que separa los mensajes en 3 categorías (urgente / útil / ignorar) con justificación de cada decisión y métrica de aciertos del modelo\\
\end{tabularx}
\normalsize

\puente{Antes de empezar, mira el viaje completo. En la sesión 5 automatizaste un trigger simple (\emph{``si X, haz Y''}). Hoy das el salto cualitativo: no solo \emph{detectar} un evento, sino \emph{interpretarlo} con IA para decidir qué hacer. La diferencia entre regla rígida y criterio aprendido es lo que distingue automatización clásica de IA aplicada.}

\begin{titlebox}{milcGradeDark}
Ruta MILC: así vamos a aprender
\end{titlebox}
\begin{tabularx}{\textwidth}{>{\centering\arraybackslash}X>{\centering\arraybackslash}X>{\centering\arraybackslash}X>{\centering\arraybackslash}X}
\phasecard{1}{milcVerde}{milcGris}{Escucha\\[-1mm]\small Observo, escucho y pregunto.} &
\phasecard{2}{milcTurquesa}{milcGris}{Sistematización\\[-1mm]\small Pensamiento computacional.} &
\phasecard{3}{milcMagenta}{milcGris}{Praxis\\[-1mm]\small Construyo y aplico.} &
\phasecard{4}{milcVino}{milcGris}{Evaluación\\[-1mm]\small Reflexión liberadora.}
\end{tabularx}


\puente{Antes de teorizar, vas a clasificar 10 mensajes reales (o simulados con bandeja de entrada típica) a mano, sin IA. Esa línea base es la referencia con la que evaluarás al modelo después. Sin línea base humana, no se puede medir si la IA es buena o mala.}

\begin{titlebox}{milcVerde}
Fase 1 · Escucha
\end{titlebox}

\begin{softbox}{milcVerde}{milcGris}{Escena para escuchar}
\textbf{Actividad 1 · IDENTIFICA --- Clasifica 10 mensajes a mano} (15 min · individual). El docente entrega 10 mensajes reales o simulados de una bandeja típica (chat de WhatsApp de negocio, correo de atención al cliente, comentarios de redes). Clasifica cada uno en una de 3 categorías: \emph{urgente} (requiere acción en menos de 2 horas), \emph{útil} (acción en el día), \emph{ignorar} (no aporta o es spam). Para cada uno, anota 1 frase de justificación. Esta es tu línea base humana, sin IA.
\end{softbox}

\checkbox Clasifiqué los 10 mensajes en una de las 3 categorías.\\
\checkbox Cada mensaje tiene justificación de 1 frase.\\
\checkbox Anoté qué mensajes me costaron más decidir (zona gris).

\begin{infoband}{milcVerde}


\textbf{Lo que va al cuaderno (Actividad 1 · IDENTIFICA).} Título: «Actividad 1 --- Línea base humana». \textbf{Formato:} tabla 10 filas × 4 columnas (Mensaje (resumen)~|~Categoría~|~Justificación~|~Dudoso?). \textbf{Sabes que terminaste cuando} reconoces qué mensajes son claros y cuáles caen en zona gris donde tu propio criterio dudó.
\end{infoband}

\puente{Ya tienes tu clasificación humana. Ahora vas a entender la \textbf{anatomía} del prompt afinado: los 5 elementos que separan un prompt amateur de uno profesional, los errores típicos que producen alucinaciones o clasificaciones inconsistentes, los principios éticos no negociables (\emph{transparencia, supervisión, derecho a corrección}).}

\begin{titlebox}{milcTurquesa}
Fase 2 · Sistematización con pensamiento computacional
\end{titlebox}

\begin{softbox}{milcTurquesa}{milcGris}{Cuatro pilares aplicados al tema}
Un prompt amateur produce clasificaciones inconsistentes. Un prompt afinado produce clasificaciones reproducibles. Vamos a descomponer los 5 elementos con los pilares del pensamiento computacional para que tu prompt tenga rigor profesional.
\end{softbox}

\small
\begin{tabularx}{\textwidth}{>{\bfseries\RaggedRight\arraybackslash}p{3.6cm}Y}
\rowcolor{milcTurquesa!90}\color{white}Pilar & \color{white}\bfseries Aplicación al tema\\
1. Descomposición & El prompt afinado se descompone en 5 elementos: (1) \textbf{Rol}: \emph{``Eres un clasificador de mensajes de atención al cliente de una panadería del barrio''}; (2) \textbf{Categorías}: exhaustivas (cubren todos los casos) y excluyentes (cada mensaje cae en una sola). \emph{urgente / útil / ignorar}; (3) \textbf{Criterios con ejemplos}: para cada categoría, 2-3 ejemplos reales: \emph{urgente: ``llegué y la panadería está cerrada''; útil: ``¿hacen sin azúcar?''; ignorar: ``Hola \emph{emoji}''}; (4) \textbf{Formato de salida}: estructurado y máquina-legible (\emph{``responde solo en JSON: \texttt{\{categoria: ``...'', justificacion: ``...''\}}''}); (5) \textbf{Borde de duda}: qué hacer cuando el modelo no está seguro (\emph{``si dudas entre urgente y útil, marca útil y agrega flag\_revision: true''}).\\
2. Reconocimiento de patrones & Todo prompt sigue el patrón \textbf{``categorías exhaustivas + criterios precisos + ejemplos + borde de duda''}: si falta alguno, el modelo improvisa. Categorías incompletas producen mensajes \emph{``perdidos''}; criterios vagos producen inconsistencia; falta de ejemplos produce alucinaciones; falta de borde de duda produce sobreconfianza falsa.\\
3. Abstracción & Lo esencial cabe en una regla: \emph{el modelo nunca clasifica mejor de lo que el humano define}. Si tu prompt es ambiguo, las salidas serán ambiguas. Si tu prompt es ético, las salidas pueden ser éticas (pero hay que supervisar). Si tu prompt es discriminatorio, las salidas serán discriminatorias. La responsabilidad ética nunca se delega al modelo.\\
4. Algoritmo conceptual & Algoritmo para afinar un prompt: (1) escribe primera versión con los 5 elementos; (2) pásalo por el modelo con 10 mensajes; (3) compara con tu línea base humana; (4) identifica los errores (falsos urgentes, falsos ignorables); (5) ajusta criterios y ejemplos del prompt; (6) repite hasta lograr 80\%+ de acierto en la muestra; (7) documenta la versión final con número de iteración.\\
\end{tabularx}
\normalsize

\begin{drawbox}{milcTurquesa}{Anatomía de un prompt afinado para clasificación}
\textbf{Rol:} contexto + autoridad del modelo. \\[1mm] \textbf{Categorías:} exhaustivas y excluyentes, con definición operativa. \\[1mm] \textbf{Criterios con ejemplos:} 2-3 ejemplos por categoría. \\[1mm] \textbf{Formato de salida:} estructurado (JSON o etiqueta+justificación). \\[1mm] \textbf{Borde de duda:} qué hacer cuando el modelo no está seguro. \\[1mm] \textbf{Restricciones éticas:} qué NO clasificar (datos sensibles, no estereotipar). \\[1mm] \textbf{Versión:} fecha + iteración del prompt.
\end{drawbox}

\begin{softbox}{milcMostaza}{milcGris}{Errores comunes y su corrección}
(1) Prompt vago: \emph{``clasifica el mensaje''} sin categorías ni criterios. (2) Categorías solapadas: \emph{urgente / importante} confunden al modelo. (3) Sin ejemplos: el modelo improvisa y produce alucinaciones. (4) Sin borde de duda: el modelo \emph{``decide''} con falsa confianza. (5) Sin supervisión humana: nadie revisa los falsos positivos. (6) Asumir que el modelo entiende contexto cultural: \emph{``listo''} significa cosas distintas en Colombia y en España.
\end{softbox}

\begin{infoband}{milcTurquesa}


\textbf{Lo que va al cuaderno (Actividad 2 · ANALIZA).} Título: «Actividad 2 --- Anatomía del prompt». \textbf{Formato:} ficha con los 5 elementos del prompt aplicados a tu caso + 5 errores comunes con marca 🚨 del que más temes + 1 línea sobre el riesgo ético principal de tu sistema. \textbf{Sabes que terminaste cuando} puedes corregir un prompt ajeno señalando vocabulario técnico (\emph{``aquí faltan categorías excluyentes''}, \emph{``aquí no hay borde de duda''}).
\end{infoband}

\puente{Tienes el método. Toca aplicarlo: redactar un prompt afinado para clasificar 20+ mensajes, ejecutarlo (ChatGPT / Claude / Gemini gratuito o en Apps Script), comparar con tu clasificación manual, medir la tasa de acierto y proponer 2 ajustes al prompt.}

\begin{titlebox}{milcMagenta}
Fase 3 · Praxis con pensamiento computacional
\end{titlebox}

\begin{softbox}{milcMagenta}{milcGris}{Construyendo el producto con los cuatro pilares}
\textbf{Actividad 3 · APLICA + EVALÚA --- Sistema de clasificación con métrica} (30 min · individual). Hora de aplicar. Redacta el prompt afinado, ejecútalo sobre 20+ mensajes reales, compara con tu línea base humana, calcula la tasa de acierto y propón 2 ajustes al prompt.
\end{softbox}

\small
\begin{tabularx}{\textwidth}{>{\bfseries\RaggedRight\arraybackslash}p{3.6cm}Y}
\rowcolor{milcMagenta!90}\color{white}Pilar & \color{white}\bfseries Aplicación al producto\\
Descomposición & Tu sistema se descompone en 5 piezas: (1) 20+ mensajes reales o simulados de tu contexto; (2) clasificación humana hecha por ti (línea base); (3) prompt afinado con los 5 elementos; (4) clasificación del modelo aplicada al mismo lote; (5) tabla comparativa con métrica de acierto + 2 ajustes propuestos.\\
Patrones & Sigue el patrón \textbf{``mide antes de creer''}: nunca asumas que el modelo clasifica bien sin medir. La tasa de acierto sobre línea base humana es el único dato confiable para decidir si el sistema sirve. Sin métrica, todo es opinión.\\
Abstracción & Cada ajuste al prompt se valida con re-ejecución: cambias el prompt, vuelves a ejecutar, re-mides. Si la métrica sube, el ajuste sirvió; si baja, revierte. La afinación es iterativa, no creativa: \emph{ciclo prompt $\to$ ejecutar $\to$ medir $\to$ ajustar}.\\
Algoritmo de armado & Algoritmo: (1) prepara tabla con 20+ mensajes; (2) clasifica a mano (línea base); (3) escribe el prompt v1; (4) pásalo por el modelo (ChatGPT, Claude, Gemini --- en chat o vía API); (5) compara y calcula \% de acierto; (6) identifica los errores; (7) ajusta prompt (v2); (8) re-ejecuta y re-mide; (9) escribe 2 propuestas concretas de mejora.\\
\end{tabularx}
\normalsize

\begin{drawbox}{milcMagenta}{Checklist antes de entregar el sistema}
\checkbox 20+ mensajes reales o simulados en tabla. \\[1mm] \checkbox Línea base humana completa con justificación. \\[1mm] \checkbox Prompt con los 5 elementos del oficio. \\[1mm] \checkbox Clasificación del modelo capturada. \\[1mm] \checkbox Métrica de acierto calculada (\% sobre línea base). \\[1mm] \checkbox 2 ajustes concretos al prompt con razón. \\[1mm] \checkbox Riesgo ético principal identificado y nombrado.
\end{drawbox}

\begin{softbox}{milcMagenta}{milcGris}{Plantilla de trabajo}
\textbf{Caso.} \textit{Contexto + dueño del proceso.} \\[1mm] \textbf{Mensajes.} \textit{20+ con resumen.} \\[1mm] \textbf{Prompt v1.} \textit{Texto completo del prompt.} \\[1mm] \textbf{Resultado.} \textit{Tabla comparativa humano vs IA + \% acierto.} \\[1mm] \textbf{Ajustes propuestos.} \textit{2 mejoras al prompt con razón.}
\end{softbox}

\begin{infoband}{milcMagenta}


\textbf{Lo que va al cuaderno (Actividad 3 · APLICA + EVALÚA).} Título: «Actividad 3 --- Sistema de clasificación». \textbf{Formato:} ficha con el prompt completo + tabla 20+ filas con humano vs modelo + métrica final + 2 ajustes propuestos. \textbf{Extensión:} 1-2 páginas de cuaderno. \textbf{Sabes que terminaste cuando} puedes decir con número si tu sistema está listo para producción o necesita más afinación.
\end{infoband}

\puente{Lo que sigue resume \emph{qué} entregas hoy: 1 prompt afinado + 1 hoja con 20+ mensajes clasificados por humano y por IA + métrica de aciertos + 2 ajustes propuestos. El criterio principal: que la métrica permita \textbf{decidir} si el sistema sirve para producción o no.}

\begin{titlebox}{milcMostaza}
Producto de la sesión
\end{titlebox}
\textbf{1 sistema de clasificación con prompt afinado + métrica de acierto}.
\begin{softbox}{milcMostaza}{milcGris}{Criterios de calidad}
(1) 20+ mensajes con clasificación humana (línea base). (2) Prompt con los 5 elementos. (3) Clasificación del modelo aplicada al mismo lote. (4) Métrica de acierto calculada. (5) 2 ajustes concretos propuestos con razón.
\end{softbox}

\puente{Antes de cerrar, mira el sistema desde las cinco dimensiones humanas. Clasificar bien con IA libera tiempo humano; clasificar mal a escala produce daño silencioso a personas reales que nunca verás.}

\begin{titlebox}{milcVino}
Fase 4 · Evaluación liberadora — cinco dimensiones
\end{titlebox}

\begin{softbox}{milcVino}{milcGris}{Sentido de la evaluación}
Evaluar no es solo revisar una nota. Es reconocer qué aprendí, qué cambió en mí, qué emociones manejé, cómo me relaciono con la ciudad y la comunidad, y qué le dejaré a quien viene detrás.
\end{softbox}

\textbf{Mi reflexión liberadora en cinco dimensiones.} Completa cada frase iniciada:

\small
\begin{tabularx}{\textwidth}{>{\bfseries\RaggedRight\arraybackslash}p{3.4cm}Y>{\RaggedRight\arraybackslash}p{4.6cm}}
\rowcolor{milcVino}\color{white}Dimensión & \color{white}\bfseries Mi respuesta (frame starter) & \color{white}\bfseries Algo que descubrí\\
1. Personal & Lo que más cambió en mí fue \linead{6.5cm} & \linead{4.2cm}\\
2. Emocional & La emoción más fuerte fue \linead{2.5cm} y la regulé \linead{3cm} & \linead{4.2cm}\\
3. Ciudadana & En democracia esto importa porque \linead{6cm} & \linead{4.2cm}\\
4. Local (Cartago) & En mi barrio sería útil para \linead{6.5cm} & \linead{4.2cm}\\
5. Intergeneracional & A mi abuela/abuelo le diría \linead{6.5cm} & \linead{4.2cm}\\
\end{tabularx}
\normalsize

\begin{infoband}{milcVino}
\textbf{Para la próxima sustentación.} Vuelve a esta tabla en seis meses. Verás que las respuestas cambian — no porque lo aprendido cambie, sino porque tú cambias. La evaluación liberadora es una conversación contigo a través del tiempo.
\end{infoband}


\puente{Y para terminar, tres voces sobre la IA en mensajería. Dussel sobre los mensajes invisibilizados por categorías, Séneca sobre la prudencia del juicio, Floridi sobre la IA como infraestructura ética de la comunicación digital.}

\begin{titlebox}{milcGradeDark}
Triángulo de pensamiento · cierre filosófico
\end{titlebox}

\begin{softbox}{milcVino}{milcGris}{Dussel · lente del nosotros}
\textit{``Toda categoría automatizada decide qué personas y qué mensajes son visibles para la institución.''}

Una IA que clasifica reclamos como \emph{``ignorar''} en cierta jerga puede silenciar a las víctimas que no escriben con el lenguaje formal esperado. Una IA que marca como \emph{``urgente''} solo lo que viene de cierto código postal puede priorizar barrios ricos sin querer. La phronesis del clasificador exige preguntar: \emph{¿qué tipo de mensaje y qué tipo de persona queda invisibilizado por mi categoría ``ignorar''?}. La responsabilidad ética es del humano que define el prompt, no del modelo que lo ejecuta.

\textbf{En mi sistema, ¿quién quedaría silenciado por la categoría \emph{ignorar}? ¿Qué cambia si reviso manualmente al menos 1 vez por semana lo que el modelo marcó como ignorable?}
\end{softbox}
\linead{\linewidth}\\[1mm]
\linead{\linewidth}

\begin{softbox}{milcMostaza}{milcGris}{Estoicismo · Séneca · cuidado interior}
\textit{``El juicio prudente exige tiempo --- la IA solo simula prudencia.''}

La IA clasifica en milisegundos, lo cual es útil para escalar pero peligroso para casos de matiz. La sabiduría estoica recomienda \emph{escalar la velocidad solo donde la velocidad no daña}. Para mensajes con matiz emocional, contexto cultural o ambigüedad, la IA debe \emph{escalar al humano} en lugar de decidir. El borde de duda en el prompt es la traducción técnica de la prudencia.

\textbf{¿Mi prompt tiene borde de duda real o asume que el modelo siempre puede decidir? ¿Qué tipo de mensaje merece siempre escalar a humano?}
\end{softbox}
\linead{\linewidth}\\[1mm]
\linead{\linewidth}

\begin{softbox}{milcTurquesa}{milcGris}{Floridi · lente de la infoesfera}
\textit{``La IA aplicada a la mensajería es infraestructura ética de la comunicación digital contemporánea.''}

Bancos, hospitales, gobiernos, plataformas usan IA para filtrar millones de mensajes diarios. Las decisiones de qué categoría asignar afectan a personas reales: deniegan créditos, retrasan diagnósticos, archivan reclamos. Aprender a diseñar estos sistemas con ética a los 16 años es alfabetización cívica decisiva: el lenguaje con el que se ejerce el poder digital sobre millones de personas.

\textbf{¿Cuántos sistemas de clasificación IA me afectaron esta semana sin que yo los viera? ¿Qué cambia si yo soy quien los diseña con criterios éticos explícitos?}
\end{softbox}
\linead{\linewidth}\\[1mm]
\linead{\linewidth}

\begin{titlebox}{milcVino}
Compromiso final
\end{titlebox}

\small
\begin{tabularx}{\textwidth}{>{\RaggedRight\arraybackslash}p{3.0cm}YYY}
\rowcolor{milcVino}\color{white}\bfseries Criterio & \color{white}\bfseries Lo logré & \color{white}\bfseries En proceso & \color{white}\bfseries Necesito apoyo\\
Comprensión del tema & Explico ideas centrales con mis palabras. & Reconozco ideas pero falta precisión. & Necesito repasar conceptos base.\\
Pensamiento computacional & Apliqué los 4 pilares al diseñar mi producto. & Apliqué algunos pilares. & Necesito releer la fase 2.\\
Aplicación y producto & Mi producto cumple los criterios y se puede revisar. & Producto funcional con ajustes pendientes. & Aún en construcción.\\
Reflexión 5 dimensiones & Respondí cada dimensión con honestidad. & Respondí algunas con detalle. & Mis respuestas son muy generales.\\
Triángulo de pensamiento & Conecté las 3 voces con mi trabajo real. & Conecté al menos una. & No relaciono las voces con mi proceso.\\
\end{tabularx}
\normalsize

\textbf{Hoy entendí que:}\\[1mm]
\linead{\linewidth}\\[1mm]
\linead{\linewidth}

\textbf{Mi compromiso para la próxima sesión es:}\\[1mm]
\linead{\linewidth}\\[1mm]
\linead{\linewidth}

\begin{softbox}{milcMostaza}{milcGris}{Cita de despedida · Marco Aurelio}
\textit{``El obstáculo en el camino se vuelve el camino. Lo que estorba al camino se vuelve camino.''}\\[1mm]
Cada error de hoy, bien observado, se convierte en pieza del camino que pisarás mañana.
\end{softbox}

\small
\begin{tabularx}{\textwidth}{>{\bfseries}p{3.6cm}Y}
\rowcolor{milcGradeDark!12}Nombre completo: & \linead{10cm}\\
Fecha: & \linead{4cm} \quad Curso/grupo: \linead{4cm}\\
Mi firma: & \linead{9cm}\\
\end{tabularx}
\normalsize

\begin{infoband}{milcGradeDark}
\textbf{Para la próxima semana.} Comparte tu sistema con un docente o emprendedor real que reciba mensajes diariamente; pide retroalimentación sobre falsos positivos o casos invisibilizados, y ajusta el prompt. La phronesis de la IA aplicada se entrena con datos reales y supervisión continua.
\end{infoband}

\end{document}
